% Industry resume, Grover Aruquipa
% Compile with pdfLaTeX (also works on Overleaf). One column on purpose:
% applicant tracking systems read single-column text far more reliably than sidebars.
\documentclass[10pt,letterpaper]{article}
\usepackage[T1]{fontenc}
\usepackage[utf8]{inputenc}
\usepackage[scaled=0.95]{helvet}
\renewcommand{\familydefault}{\sfdefault}
\usepackage[top=1.0cm,bottom=0.8cm,left=1.5cm,right=1.5cm]{geometry}
\usepackage{enumitem}
\usepackage{titlesec}
\usepackage{xcolor}
\usepackage{microtype}
\usepackage[hidelinks]{hyperref}

\definecolor{ink}{HTML}{111111}
\definecolor{accent}{HTML}{0D2B61}
\definecolor{muted}{HTML}{555555}
\color{ink}
\pagestyle{empty}
\setlength{\parindent}{0pt}
\setlength{\parskip}{0pt}

\titleformat{\section}{\normalsize\bfseries\color{accent}}{}{0em}{\MakeUppercase}[\vspace{-6pt}{\color{accent}\rule{\linewidth}{0.6pt}}]
\titlespacing*{\section}{0pt}{7pt}{3pt}
\setlist[itemize]{leftmargin=1.1em,itemsep=0.5pt,topsep=1pt,parsep=0pt,label=\textbullet}

% \role{Title}{Organization}{Location}{Dates}
\newcommand{\role}[4]{\textbf{#1}, #2 \hfill {\color{muted}#3 \textbar\ #4}\par}

\begin{document}

%================ HEADER ================
{\LARGE\bfseries Grover Aruquipa}\par\vspace{2pt}
{\large Mechanical Design Engineer, Robotic Hand \textbar\ Linkages and Grippers, CAD to Tested Prototype}\par\vspace{3pt}
{\small Qu\'ebec, QC, Canada \ \textbullet\ \
\href{mailto:grover.aruquipa.9@gmail.com}{grover.aruquipa.9@gmail.com} \ \textbullet\ \
\href{https://www.linkedin.com/in/grover-aruquipa-996917126}{linkedin.com/in/grover-aruquipa-996917126} \ \textbullet\ \
\href{https://groveraruquipa.github.io}{groveraruquipa.github.io}}\par\vspace{1pt}
{\small\bfseries Available from January 2027 \ \textbullet\ \ Open to relocation}\par


%================ SUMMARY ================
\section{Summary}
Mechanical design engineer finishing a PhD in robot mechanisms and grasping (Universit\'e Laval, defense early 2027),
author of a published review of parallel mechanisms for grasping and robotic hands. I have designed, built and tested three
robots whose grippers are driven through linkages and concentric shafts from base-mounted motors, doing the whole loop
myself: linkage synthesis, SolidWorks, material and process selection, printing and machining, integration and test.
Eligible for an open Post-Graduation Work Permit in Canada.

%================ SKILLS ================
\section{Technical Skills}
\begin{itemize}
\item \textbf{Mechanical design:} SolidWorks (expert, three full robot assemblies), Creo; robotic linkages and mechanisms
(Sarrus, parallelogram, spherical and parallel linkages), joints, actuators and transmissions (hollow shafts,
belts, gears), integrated grippers and fingers; drawings and BOMs; material and process selection for
strength, weight and cost.
\item \textbf{Analysis:} statics, dynamics and kinematics of linkages, grasp and force analysis, Jacobians and singularities,
FEA (ANSYS); simulation in MATLAB, Simscape, MuJoCo, NVIDIA Isaac Sim and Pinocchio.
\item \textbf{Prototyping and test:} 3D printing (SLS, FDM) and CNC machining of prototypes, assembly, motor and encoder
integration, performance and durability tests (repeatability, tracking, continuous rotation); real-time control at
500\,Hz (OPAL-RT, Simulink Real-Time).
\item \textbf{Software and electronics:} Python, C++ (real time), MATLAB; ROS and ROS~2, MoveIt, OpenCV, PyTorch, Git,
Linux; Altium (PCB), motor drives, wiring.
\end{itemize}

%================ EXPERIENCE ================
\section{Experience}
\role{Doctoral Researcher}{Robotics Laboratory, Universit\'e Laval}{Qu\'ebec, Canada}{2023 -- present}
\begin{itemize}
\item \textbf{4-DoF coaxial wrist-gripper:} four actuation channels through concentric hollow shafts, non-backdrivable
gripper, 180\,mm footprint, 0.25\,kg moving mass; unlimited roll verified over ten continuous revolutions with
sub-degree tracking while holding the object.
\item \textbf{Four-leg Delta robot with integrated Sarrus-linkage gripper (in progress):} grasp force estimated from motor
currents, without a sensor at the fingers, to pick fragile objects such as eggs quickly.
\item \textbf{(6+1)-DoF parallel robot:} seven base-mounted actuators; gripper actuated from the base through a concentric
tube while the hand rotates without limit.
\item Wrote a published review of parallel mechanisms for grasping, including parallel and underactuated robotic hands
(ASME JMR). Day-to-day lead of two graduate students on an 8-actuator robot; reviewer for ASME J.\ Mech.\ Robotics.
\end{itemize}

\vspace{3pt}
\role{Visiting Researcher}{Institute of Robotics, Johannes Kepler University Linz}{Linz, Austria}{Apr. -- May 2026}
\begin{itemize}
\item Developed singularity and grasp-quality indices for manipulators with integrated grasping, hosted by Univ.-Prof.~A.~M\"uller (Mitacs Globalink Research Award).
\end{itemize}

\vspace{3pt}
\role{Research Intern}{FEMTO-ST (CNRS)}{Besan\c{c}on, France}{2022 -- 2023}
\begin{itemize}
\item Designed and modeled a redundant dexterous parallel robot with two unlimited rotations for micro-manipulation, and wrote its real-time C++ control software; clean-room fabrication projects.
\end{itemize}

\vspace{3pt}
\textbf{Industry roles} \hfill {\color{muted}2020 -- 2021}\par
{\small Computer Vision Engineer, Sencon Srl (industrial vision, scaffolding detection) \ \textbullet\ \ Data Science
Intern, Hitachi ABB Power Grids, Zurich (machine learning for high-voltage circuit breakers) \ \textbullet\ \ Data
Scientist, KTH Royal Institute of Technology.}

%================ EDUCATION ================
\section{Education}
\textbf{PhD, Mechanical Engineering}, Universit\'e Laval \hfill {\color{muted}Qu\'ebec, Canada \textbar\ 2023 -- 2027 (expected)}\par
{\small Kinematically redundant parallel robots with unlimited rotations and integrated grasping. Advisors: C.~Gosselin, P.~Lambert.}\par\vspace{2pt}
\textbf{MSc, Smart Integrated Systems}, Universit\'e Bourgogne--Franche-Comt\'e \hfill {\color{muted}Besan\c{c}on, France \textbar\ 2021 -- 2023}\par\vspace{2pt}
\textbf{BEng, Mechatronics Engineering}, Bolivian Catholic University \hfill {\color{muted}La Paz, Bolivia \textbar\ 2014 -- 2019}\par
{\small Thesis: reconfigurable Delta robot for product picking.}

%================ AWARDS ================
\section{Awards and Languages}
\begin{itemize}
\item 1st Place, Graduate Division, ASME Student Mechanisms and Robotics Design Competition (2025), team lead.
\item Finalist, MassRobotics competition, Boston, with one of my robot prototypes, team lead.
\item Mitacs Globalink Research Award (2026); 1st Place, Student Paper Competition, Argentine Robotics Conference (2019).
\item Languages: Spanish (native), English (full professional), French (intermediate).
\end{itemize}

\newpage
%================ PUBLICATIONS
\section{Selected Publications}
{\small First author on all papers. Full list at groveraruquipa.github.io.}\par\vspace{2pt}
\begin{itemize}
\item \textbf{G.~Aruquipa}, P.~Lambert, R.~Dahmouche, C.~Gosselin, ``Parallel robots for grasping tasks: a review,'' \emph{ASME Journal of Mechanisms and Robotics}, vol.~18, no.~12, art.~120801, 2026. doi:10.1115/1.4072612
\item \textbf{G.~Aruquipa}, P.~Lambert, C.~Gosselin, ``A compact coaxial parallel wrist-gripper with unlimited tool-axis rotation,'' under review, \emph{Robotics and Autonomous Systems}.
\item \textbf{G.~Aruquipa}, P.~Lambert, C.~Gosselin, ``Analysis and design of a novel (6+1)-DoF parallel robot with unlimited rotation and grasping capabilities,'' \emph{Mechanism and Machine Theory}, vol.~224, art.~106438, 2026.
\item \textbf{G.~Aruquipa}, V.~Muralidharan, P.~Lambert, C.~Gosselin, ``Analysis and design of a novel compact four-degree-of-freedom parallel robot with unlimited rotation,'' \emph{ASME Journal of Mechanisms and Robotics}, vol.~18, no.~6, art.~061001, 2026.
\item \textbf{G.~Aruquipa}, P.~Lambert, C.~Gosselin, ``Kinematic analysis and design of a novel 9-DoF parallel robot with grasping capabilities,'' \emph{CCToMM Symposium on Mechanisms, Machines and Mechatronics}, 2025, pp.~117--127.
\item \textbf{G.~Aruquipa}, P.~Lambert, C.~Gosselin, ``A kinematically redundant (6+2)-DoF parallel robot with three orthogonal unlimited rotations and coaxial actuation,'' submitted to \emph{IEEE ICRA 2027}.
\end{itemize}

%================ REFERENCES ================
\section{References}
\begin{itemize}
\item \textbf{Prof.~Cl\'ement Gosselin}, Full Professor, Department of Mechanical Engineering, Universit\'e Laval, Qu\'ebec,
Canada. PhD supervisor. \href{mailto:clement.gosselin@gmc.ulaval.ca}{clement.gosselin@gmc.ulaval.ca}
\item \textbf{Prof.~Patrice Lambert}, Assistant Professor, Department of Mechanical Engineering, Universit\'e Laval,
Qu\'ebec, Canada. PhD co-supervisor. \href{mailto:patrice.lambert@gmc.ulaval.ca}{patrice.lambert@gmc.ulaval.ca}
\item \textbf{Univ.-Prof.~Dr.-Ing.~Andreas M\"uller}, Head of the Institute of Robotics, Johannes Kepler University Linz,
Austria. Host during my research stay (2026). \href{mailto:a.mueller@jku.at}{a.mueller@jku.at}
\end{itemize}

\end{document}
